\begingroup\fontsize{8}{10}\selectfont

\begin{longtable}[t]{lp{4.2cm}p{8.2cm}}
\caption{\label{tab:appD_quotes_part2}Open-ended responses: valuation, cost--benefit, and overall. Verbatim quotes; empty entries and pure `N/A' answers are omitted.}\\
\toprule
Section & Prompt & Response\\
\midrule
\endfirsthead
\caption[]{Open-ended responses: valuation, cost--benefit, and overall. Verbatim quotes; empty entries and pure `N/A' answers are omitted. \textit{(continued)}}\\
\toprule
Section & Prompt & Response\\
\midrule
\endhead

\endfoot
\bottomrule
\endlastfoot
Cost--benefit & Additional cost considerations to incorporate & Just be certain that the cost of a flight is not JUST the cost of that flight.  A flight can only happen because of the staff and the long-term maintenance... not just the initial cost of the plane and the fuel it uses during a flight. So more of a bulk average cost of a flight (harder to get at).\\
 &  & I believe even the additional costs for maintenance, etc. would again be cost effective vs damage costs.\\
 &  & Changes in the performance baseline given degraded observations say with the loss of DMSP and POES. Instrumentation and availability of aircraft should also be considered. There are also external pressures to replace some NOAA observations with commercial providers or drones.\\
 &  & Staffing is also a consideratioj.\\
 &  & What is missing is a comparison to the other avenues of improving forecasts, such as increased computer power for the forecasts and the use of AI forecast models.\\
\addlinespace
Cost--benefit & Further analysis to increase confidence in the policy implication & Neutral\\
 &  & Identifying which counties/populations are most impacted by hurricanes and supported by NOAA information. I know demographics and socioeconomics are deprioritized currently, but not everyone has the same experience or means to respond to hurricane information\\
 &  & Addressing the missing components to the forecast improvements.\\
 &  & Not sure.\\
Cost--benefit & Whether the benefit--cost ratio justifies fleet recapitalisation & Even if the value is off slightly because this-or-that assumption, that ratio is very high and easily justifies fleet recapitalization. Even it were 10:1!\\
\addlinespace
 &  & Absolutely. This study mostly considers the costs / financials, which is clearly justified. NOAAs mission focuses on life and property, lives saved vs lost should also be considered as value, again justifying fleet recapitalization.\\
 &  & BCRs are not used in decisionmaking as they don't consider time. Net present value is preferred. Justification is much more complex than just one summary statistic.\\
 &  & Although there are uncertainties, the benefit-to-cost ratio is very large so the  qualitatively conclusion probably wouldn't change. Also, for reasons described above, it it probably an underestimate.\\
 &  & Yes.\\
Overall & Innovative extension that could substantially increase the policy impact & I need to spend more time with the paper\\
\addlinespace
 &  & Neutral\\
 &  & Understanding the full breadth of observations and models in the problem space and identifying the percent contribution of each to the final dollar value.\\
 &  & Could the methodology be applied to past cases to compare the actual forecast uncertainty with the post-storm damage estimates?\\
 &  & Also maybe comparison to other government supported programs that provide economic benefits.\\
Overall & Single additional analysis that would most strengthen the paper & I need to spend more time with the paper\\
\addlinespace
 &  & Neutral\\
 &  & Statistical robustness, uncertainty, and limitations discussion\\
 &  & More specifically addressing the impact of track and size improvements.\\
 &  & As we discussed, normalizing the improvements in wind errors and cost benefits compared to other measurable changes, such the changes over the last 20 years.\\
Overall & Stage of the analysis presenting the greatest vulnerability & Neutral\\
\addlinespace
 &  & Stakeholder engagement, data collection, and final estimation strategy.\\
 &  & The economic part of the cost analysis is the most empirical is probably easier to criticize.\\
 &  & I still think there may be large errors with the wind swaths, and further improvements in the shape of the wind field, and the size, would lead to more accurate results.\\
Valuation & Categories of benefits or costs that may not be fully incorporated & I need to spend more time with the paper to thoroughly understand this relationship.\\
 &  & Does a good forecast of an impactful hurricane reduce losses because the right people are taking mitigating preparations?\\
\addlinespace
 &  & Lives\\
 &  & There may be policy considerations missing given the study framing. Investments in wetlands, infrastructure, modeling, or other decisionmaking may help as much or in different ways rather than augmenting observations further.\\
 &  & I am not sure if the reduction in the loss of life was incorporated.\\
 &  & Do not know.\\
Valuation & Sensitivity of monetary projections to key assumptions & I need to spend more time with the paper to thoroughly understand this relationship.\\
\addlinespace
 &  & Neutral\\
 &  & No informed enough to answer.\\
 &  & Given climate variability, hurricane frequency and severity may be more unpredictable. Given the large dollar values attributed to changes in beta, I would ensure robustness checks and sensitivity analyses.\\
 &  & I don't have a lot of experience in this area, but given the nonlinear relationships the projections should be very sensitive to the assumptions.\\
 &  & Hard for me to say.\\
\addlinespace
Valuation & Whether the valuation captures the full social value of forecasts & I need to spend more time with the paper to thoroughly understand this relationship.\\
 &  & Neutral\\
 &  & I think there are additional ways and measures to capture the full social value of improving forecast. Damage impacts (not just monetarily, but indirect examples too) / lives should be metrics as well.\\
 &  & Not the full value but a meaningful subset.\\
 &  & I think it is an under-estimate because storm size and track error reductions and impacts from storm surge and rainfall are not fully accounted for.\\
\addlinespace
 &  & As we discussed in the meeting, it's likely that the valuation underestimates the full social value.\\*
\end{longtable}
\endgroup{}
